\scanpage{9}
% archive sheet number: 8
\mspage{- IV-\handwritten{965} -}
% Two long oblique strokes cancel the lower corollary; a wavy stroke also crosses the remark.

\uline{Corollaire} (18.9.4) .- \uline{Soit} $A$ \uline{un anneau local noethérien , hensélien et dont les fibres formelles sont géométriquement normales} \deleted{\illegible{1}} (par exemple un anneau local noethérien hensélien et \uline{excellent}) , \uline{et posons} $A'=\hat A$ . \uline{Soient} $X$ \uline{un} $A$-\uline{préschéma} , $X'=X\handwritten{\otimes}_A A'$ , $g:\ X'\handwritten{\to} X$ \uline{la projection canonique . Alors l'application} $U\handwritten{\rightsquigarrow} g^{-1}(U)$ \uline{est une bijection de l'ensemble des parties à la fois ouvertes et fermées de} $X$ \uline{sur l'ensemble des parties à la fois ouvertes et fermées de} $X'$ .

Comme $g$ est un morphisme \deleted{\unsure{très|f}} fidèlement plat et quasi-compact , on sait (2.3.12) que la topologie de $X$ est \uline{quotient} de celle de $X'$ par la relation d'équivalence définie par $f$ . Tout revient donc à voir \deleted{\unsure{que}} qu'une partie ouverte et fermée $U'$ de $X'$ est \uline{saturée} pour cette relation . Or , comme le morphisme $\operatorname{Spec}(A')\handwritten{\to}\operatorname{Spec}(A)$ \deleted{\unsure{est}} a ses fibres géométriquement connexes (18.9.1) , les fibres de $g$ sont \uline{connexes} , d'où aussitôt notre assertion .

En particulier , si $X$ est localement connexe (ce qui sera le cas si $X$ est localement noethérien) , alors , pour toute \deleted{\unsure{partie connexe}} composante connexe $U$ de $X$ \handwritten{(} qui est ouverte et fermée dans $X$ \handwritten{),} $g^{-1}(U)$ est connexe (Bourbaki, \uline{Top.gén.},chap.I,$3^{\mathrm e}$ éd.,\S\ 11,n$^{\mathrm o}$3 ,\handwritten{prop 7} ) . Si l'on désigne par $\handwritten{\pi}_0(X)$ l'ensemble des composantes connexes de $X$ , , l'application $\handwritten{\pi}_0(X')\handwritten{\to}\handwritten{\pi}_0(X)$ déduite canoniquement de $g$ est donc \uline{bijective} .

\deleted{\uline{Corollaire} (18.9.5) .- \uline{Sous les hypothèses de} (18.9.4) , \uline{le foncteur}}

% slanted
\margin{\handwritten{``séparés'' inutile !}}

\deleted{\[
Z\handwritten{\rightsquigarrow} Z\handwritten{\otimes}_A A'\tag*{(18.9.5.1)}
\]}

\deleted{\uline{de la catégorie des} \deleted{$X$-}\inserted{\handwritten{pré}}\uline{schémas étales} \deleted{\handwritten{(resp. étales et finis)}} \inserted{\handwritten{et séparés}} \uline{sur} $X$ \uline{dans celle des} \deleted{\illegible{2}} \deleted{\unsure{$X'$-|$X$-}}\inserted{\handwritten{pré}}\uline{schémas étales} \deleted{\handwritten{(resp. étales et finis)}} \inserted{\handwritten{et séparés}} \uline{sur} $X'$ , \uline{est pleinement fidèle} .}

% slanted
\margin{\handwritten{la \unsure{var.|vers.|vrai} n'apporte rien ; $X$-schéma risque d'induire en erreur ---}}

\deleted{En effet , si $Z_1,Z_2$ sont deux $X$-schémas étales sur $X$ , $\operatorname{Hom}_X(Z_1,Z_2)$ s'identifie à $\handwritten{\Gamma}(Z_1\handwritten{\times}_X Z_2/Z_1)$ , et $Z_1\handwritten{\times}_X Z_2$ est étale et séparé sur $Z_1$ , donc finalement (17.9.3) , $\operatorname{Hom}_X(Z_1,Z_2)$ s'identifie fonctoriellement à l'ensemble des parties ouvertes et fermées $U$ de $Z_1\handwritten{\times}_X Z_2$ telles que la restriction de la projection $p:\ Z_1\handwritten{\times}_X Z_2\handwritten{\to} Z_1$ à $U$ soit un morphisme surjectif et radiciel . L\deleted{a}\inserted{e} corollaire résulte donc aussitôt de (18.9.4) et du fait que , si $Z_i'=Z_i\handwritten{\otimes}_A A'$ \deleted{\illegible{1}} $(i=1,2)$ , on a $Z_1'\handwritten{\times}_{X'} Z_2'=(Z_1\handwritten{\times}_X Z_2)\handwritten{\otimes}_A A'$ .}

\deleted{\uline{Remarques} (18.9.6) .- \deleted{\unsure{i)} Dans tous les cas connus} \inserted{\handwritten{Il est possible que si le corps résiduel de $A$ est de car. nulle,}} le foncteur (18.9.5.1) \deleted{\unsure{est}} \inserted{\handwritten{induit}}}
